\documentclass[12pt]{article}
\usepackage[margin=1.1in]{geometry}
\usepackage{amsmath,amssymb}
\usepackage{booktabs,multirow}
\usepackage{graphicx}
\usepackage{natbib}
\usepackage{setspace}
\usepackage[hidelinks]{hyperref}
\usepackage{caption}
\newcommand{\champForecastLL}{0.324}
\newcommand{\champForecastAcc}{86.2}
\newcommand{\champContestedAcc}{85.1}
\newcommand{\champAPRE}{0.48}
\newcommand{\incForecastLL}{0.348}
\newcommand{\champGain}{0.024}
\newcommand{\champCongressoutLL}{0.643}
\newcommand{\majorityCongressoutLL}{0.617}
\newcommand{\gbRawForecastLL}{0.803}
\newcommand{\prospectiveVoneLL}{0.331}
\newcommand{\prospectiveVoneAcc}{85.7}
\newcommand{\prospectiveVtwoLL}{0.332}
\newcommand{\prospectiveVtwoAcc}{86.8}
\newcommand{\prospectiveN}{9,451}
\newcommand{\regimeAbestLL}{0.153}

% auto-generated context
\newcommand{\ukRoll}{1235}
\newcommand{\ukDate}{2024-12-20}
\newcommand{\ukCut}{3.44}
\newcommand{\ukDiscrim}{-1.11}
\newcommand{\ukYeaShare}{92}
\newcommand{\ukNPred}{400}

\newcommand{\revoltProbeR}{0.25}
\newcommand{\revoltProbeN}{944}
\newcommand{\ssfaAUCtext}{0.74}
\newcommand{\ssfaDefectors}{71}
\newcommand{\ssfaAUCnotext}{0.69}

\newcommand{\seenLLchamp}{0.504}
\newcommand{\freshLLchamp}{0.300}
\newcommand{\seenLLnotext}{0.590}
\newcommand{\freshLLnotext}{0.421}
\newcommand{\seenLLclean}{0.518}
\newcommand{\freshLLclean}{0.325}
\newcommand{\seenRC}{368}
\newcommand{\freshRC}{1,911}
\newcommand{\houseLLchamp}{0.306}
\newcommand{\senateLLchamp}{0.464}
\newcommand{\houseLLnotext}{0.434}
\newcommand{\senateLLnotext}{0.480}

\newcommand{\ablOrigLL}{0.328}
\newcommand{\ablOrigAUC}{0.78}
\newcommand{\ablShuffLL}{0.392}
\newcommand{\ablShuffAUC}{0.78}
\newcommand{\ablNoSumLL}{0.351}
\newcommand{\ablNoQLL}{0.383}
\newcommand{\ablNoPropLL}{0.339}
\newcommand{\ablGenTitleLL}{0.340}
\newcommand{\ablRC}{1,957}

\newcommand{\brierChamp}{0.099}
\newcommand{\eceChamp}{0.017}
\newcommand{\defShareRChamp}{0.34}
\newcommand{\brierClean}{0.107}
\newcommand{\eceClean}{0.021}
\newcommand{\defShareRClean}{0.32}
\newcommand{\brierNotext}{0.147}
\newcommand{\eceNotext}{0.028}
\newcommand{\defShareRNotext}{0.27}
\newcommand{\defShareNrc}{995}

\newcommand{\overlapRho}{0.05}
\newcommand{\overlapN}{659}
\newcommand{\paulDefGainPct}{77}
\newcommand{\amashDefGainPct}{73}

\newcommand{\cfAraOrig}{6.5}
\newcommand{\cfAraGeneric}{2.7}
\newcommand{\cfAraPassage}{13.5}
\newcommand{\cfAraBio}{8.4}
\newcommand{\cfArdaOrig}{14.5}
\newcommand{\cfArdaCut}{14.7}
\newcommand{\cfArdaGeneric}{12.6}
\newcommand{\cfArdaSwap}{7.3}

\newcommand{\cutHiDiscR}{0.70}
\newcommand{\cutHiDiscMAE}{0.35}
\newcommand{\cutPassageR}{0.61}
\newcommand{\cutPassageMAE}{0.82}

\newcommand{\paulHeldoutN}{949}
\newcommand{\amashHeldoutN}{579}


\title{Reading the Bills: Text, Votes, and the\\ Measurement of Congress}
\author{Andrew B.\ Hall\thanks{Stanford University. Replication materials,
including the full pipeline, model registry, pre-registered experiment
ledger, and session logs of the AI-assisted development process, are
available in the project repository. Every statistic and figure is
generated programmatically from the analysis code's output files; no
number is transcribed by hand.}}
\date{Draft: July 2026 --- preliminary, comments welcome}

\begin{document}
\maketitle
\onehalfspacing

\begin{abstract}
\noindent The models political science uses to measure Congress
predict rollcall votes without reading the bills being voted on. We
construct an open benchmark of 10.6 million individual voting
decisions (101st--119th Congresses) and evaluate an ensemble whose
language-model encoders take the text of each measure as input. Three
results follow. First, the approach outperforms existing ones: on
identical held-out votes it improves on DW-NOMINATE under the
classification literature's own fit statistic, with the largest gains
concentrated among the members one-dimensional scaling is known to
mis-measure (Ron Paul's geometric mean probability rises from
\paulGMPnom{} to \paulGMPours), and on strictly future votes a frozen,
hash-pinned artifact sustains \champForecastAcc\% accuracy---beyond
both classical spatial models and the earlier text-based literature.
Second, the text is what allows the model to anticipate protest
voting: across \protRC{} held-out majority-side rollcalls, reading
the bill raises within-vote defection detection from \protAUCnotext{}
to \protAUCtext{} AUC, and embedding-space diagnostics show how---the
encoder places bills that provoke intra-party revolts near earlier
bills that provoked them, and a linear probe on the text alone
predicts each rollcall's majority-defection share out of sample.
Third, text locates bills before anyone votes: on \cutTestN{}
held-out future rollcalls, text plus pre-vote metadata predicts
realized cutpoints at $r = \cutComboR$ and coalition direction at
\dirCombo\%, a prospective counterpart to a bill-location literature
that has required realized votes. We release the benchmark, the
models, an interactive forecaster, and a pre-registered ledger of
every experiment, positive and negative.
\end{abstract}

\section{Introduction}

Nearly everything political science measures about Congress---ideal
points, party unity, issue positions---derives from models of
rollcall votes that treat the bill under consideration as an indexed
event with estimable parameters rather than as a text with content
\citep{poole1997congress, clinton2004statistical}. A literature at
the intersection of machine learning and legislative studies sought
to change this with the text tools of its era
\citep{gerrish2011predicting, kraft2016embedding}, and the
measurement literature has since pressed on the classical model from
several directions: unpredictable voters
\citep{lauderdale2010unpredictable}, issue-specific dimensions
\citep{lauderdale2014scaling, crespin2010dimensions}, estimators for
massive vote matrices \citep{imai2016fast}, machine learning as a
measurement strategy \citep{bonica2018inferring}, and, most recently,
purpose-built models for the members whom one-dimensional scaling
visibly mis-measures---extremists who vote with the opposite pole
\citep{duckmayr2023ends} and majority members who cast protest votes
against bills they arguably prefer \citep{fowler2026protest}. Modern
language models read legislative text far better than anything
available to the earlier program. This paper asks what that capacity
contributes to the measurement of Congress and offers a three-part
answer.

First, the approach outperforms existing ones
(Section~\ref{sec:outperform}). On an open benchmark of 10.6 million
voting decisions spanning the 101st through 119th Congresses, models
that read the bills improve on the classical toolkit in both of the
settings that toolkit defines. In the psychometric completion
setting---a vote-only exercise in which no text enters---an
expressive spatial model fits held-out votes better than frozen
DW-NOMINATE for \memberGMPshare\% of members on the NOMINATE
literature's own statistic, with the improvements concentrated in
precisely the cases that motivated a decade of bespoke scaling
models. In genuine
forecasting, a text-reading ensemble attains \champForecastAcc\%
accuracy on strictly future votes (\champContestedAcc\% on contested
votes); the 2011--2016 text-based models, reimplemented and evaluated
under identical conditions, do not improve on a simple table of
member voting rates. A frozen artifact, hash-pinned before the votes
it is scored on were cast, verifies that these figures do not reflect
researcher iteration on the test set.

Second, reading the bill is what allows the model to anticipate
protest voting (Section~\ref{sec:protest}). Protest voting---majority
members defecting on bills they arguably prefer---is the newest
concern of the scaling literature \citep{fowler2026protest,
duckmayr2023ends}, where it has been treated retrospectively, by
reinterpreting votes already cast. We show that protest voting is
predictable in advance and locate the source of the predictability in
the text. Member history identifies who sits on a bill's losing
flank; only the bill's content indicates which bills put the flank in
revolt. We document the mechanism directly: in the encoder's
embedding space, a shutdown-averting continuing resolution sits in a
neighborhood of earlier must-pass fiscal packages that provoked the
same right-flank defections, and a linear probe on the text alone
predicts each holdout rollcall's majority-defection share with no
member-level information whatsoever.

Third, text predicts where a bill will divide the chamber before
anyone votes (Section~\ref{sec:cutpoints}). The bill-location
literature estimates proposal and status-quo positions from realized
votes and cosponsorships \citep{peress2013estimating,
richman2011parties}; by construction it cannot locate a bill that has
not been voted on. We show that a bill's text can: on held-out future
rollcalls, text plus pre-vote metadata predicts realized cutpoint
locations at $r = \cutComboR$ and coalition direction at \dirCombo\%.
The division of labor among the inputs is itself informative:
institutional metadata indicates where a vote will cut, the acting
legislator's party supplies a default side, and the bill's content
indicates when the default is wrong. An interactive forecaster
operationalizes the result for hypothetical bills.

The remainder of what this exercise produced---evaluation-regime
audits, negative results with diagnosed mechanisms, the behavior of
member history across majority transitions---appears in the
appendices and repository. Section~\ref{sec:limits} states the two
boundaries a user of these methods should know.

\section{Benchmark, Models, and Definitions}
\label{sec:setup}

\textbf{Task and data.} We predict individual voting decisions: given
everything observable before a rollcall, each member's probability of
voting yea, conditional on a recorded yea or nay. The panel comprises
every such decision in the 101st--119th Congresses from Voteview
\citep{lewis2026voteview}; bill-side information (titles, dated
summary versions, sponsors, amendment records) comes from government
BILLSTATUS files, which begin with the 108th Congress. All
forecasting results therefore use the text-era congresses
(108th--119th); earlier congresses enter only the no-text completion
fits. Splits are deterministic and content-hashed, the evaluation
harness strips labels before models see evaluation rows, and the
feature-timing audit appears in Appendix~\ref{app:audit}.

\textbf{Evaluation.} We evaluate under progressively stricter
information constraints: \emph{completion}, in which random votes are
held out and same-rollcall votes remain observable (the psychometric
setting in which ideal-point models are traditionally assessed);
\emph{temporal forecast}, in which each congress-chamber's final ten
percent of rollcalls by date is held out; and \emph{prospective}, in
which hash-pinned frozen artifacts are scored only on votes cast
after they were finalized, following the reusable-holdout logic of
\citet{dwork2015holdout}. Metrics are the field's standard ones:
correct classification, APRE, and geometric mean probability
\citep{carroll2013structure}, with log loss ($= -\log \mathrm{GMP}$)
as the primary scoring rule.

\textbf{Models.} The forecasting model is a logistic blend of
two-tower networks: member embeddings interact with rollcall
parameters that a deep head amortizes from pre-vote text (the vote
question, the vote description, and the latest bill summary dated
before the vote), anchored on member-by-question empirical rates and
calibrated on a temporal slice of the training data. Its no-text
counterpart---the identical architecture with all text inputs
deleted---serves as the counterfactual throughout. The classical
comparators are logistic ideal-point models, a frozen-DW-NOMINATE
model, and reimplementations of the 2011--2016 text-based models.
The full specification appears in Appendix~\ref{app:spec};
development followed a pre-registered experiment ledger
(Appendix~\ref{app:ledger}).

\textbf{Definitions.} The spatial fits estimate $P(\text{yea}) =
\sigma(a_v x_i + b_v + c_i)$ per chamber-congress (party-sign
initialized); a rollcall's \emph{average-member cutpoint} is
$-(b_v + \bar c)/a_v$, reported when $|a_v| \geq 0.35$ on the fitted
scale, with \emph{direction} $\mathrm{sign}(a_v)$. All derived
quantities are conditional on the realized rollcall agenda. The
identifying normalization: member positions are standardized within
chamber-congress (mean zero, unit variance) and oriented so that the
Republican mean exceeds the Democratic mean; cutpoints, the
$|a_v|$ screen, and directions are stated on that scale.

\section{The Approach Outperforms Existing Ones}
\label{sec:outperform}

The classical toolkit is assessed in two settings---completing a vote
matrix and forecasting future votes---and we take each in turn.

\subsection{Held-out fit against the canonical measure}

The completion regime involves no text: it asks only how well a
spatial model of votes fits held-out votes, and so isolates model
flexibility from bill content. In it, an eight-dimensional spatial
model reaches a geometric mean probability of \regimeAbestGMP{}
(\regimeAbestAcc\% correct classification) on held-out votes, against
\regimeAnominateGMP{} for frozen DW-NOMINATE positions with estimated
rollcall parameters; the full ladder of intermediate models appears
in Appendix Table~\ref{tab:regimeA}. Both models are fit on identical
training cells and scored on held-out cells that neither model's
parameters---including the rollcall parameters---were estimated on.
Figure~\ref{fig:membergmp} disaggregates the comparison to the member
level, computing each member's held-out GMP under both models on
identical votes. Two
features of the figure organize the result. For the typical member,
DW-NOMINATE is nearly unimprovable: the median member GMP is
\memberGMPmedNom{} under the classical model and \memberGMPmedOurs{}
under ours, and the richer model is higher for \memberGMPshare\% of
members. The improvements instead concentrate in the tail---and the
tail is composed of exactly the members for whom the past decade of
scaling research built bespoke models. Ron Paul's held-out GMP rises
from \paulGMPnom{} to \paulGMPours. The other labeled members---
Justin Amash, Helen Chenoweth-Hage, Mark Neumann, Timothy
Penny---are ends-against-the-middle voters and habitual defectors of
the kind modeled by \citet{duckmayr2023ends} and
\citet{fowler2026protest}; here a single flexible model absorbs them
without member-specific machinery. Table~\ref{tab:hardcase} takes the
two canonical cases vote by vote: the held-out votes on which the
flexible model most improves are exactly the votes that made these
members famous cases. For Paul (\paulHeldoutN{} held-out votes), the
largest improvements are nay votes on defense authorizations and
appropriations, an intelligence authorization, the Electronic
Surveillance Modernization Act, and an Iran-focused resolution---votes
the one-dimensional model prices below two percent and the flexible
model above ninety-five. For Amash (\amashHeldoutN{}), they are rules
resolutions, a commemorative commission extension, an
Israeli--Palestinian policy resolution, and the FY2019 national
defense authorization. The improvements, in short, sit on the
libertarian and procedural-objector votes the bespoke-model
literature describes, recovered here without member-specific
machinery. Convergent validity is preserved:
the model's first-dimension member positions (the oriented leading
coordinate of the fitted space) correlate with first-dimension
DW-NOMINATE at $r = .98$, and per-congress with Nokken--Poole scores
\citep{nokken2004congressional} at $r = \npRmin$--$\npRmax$, so the
improvement in fit does not come at the cost of the scale's
interpretation.

\begin{table}[t]
\centering
\caption{The canonical hard cases, vote by vote: held-out votes with
the largest improvement in the probability assigned to the vote the
member actually cast, flexible spatial model versus the
frozen-DW-NOMINATE model (identical held-out cells).}
\label{tab:hardcase}
{\small \begin{tabular}{p{7.2cm}cccc}
\toprule
 & & & \multicolumn{2}{c}{$P(\text{vote cast})$} \\
\cmidrule(lr){4-5}
Held-out vote & Congress & Cast & DW-NOM & This paper \\
\midrule
\multicolumn{5}{l}{\emph{Ron Paul (R-TX)}: 949 held-out votes} \\
\quad National Defense Authorization Act for Fisca\dots & 109 & Nay & 0.00 & 0.99 \\
\quad Intelligence Authorization, FY 1998 & 105 & Nay & 0.00 & 0.99 \\
\quad Electronic Surveillance Modernization Act & 109 & Nay & 0.00 & 0.97 \\
\quad Department of Defense Appropriations Act for\dots & 106 & Nay & 0.02 & 0.99 \\
\quad A concurrent resolution expressing the conce\dots & 108 & Nay & 0.00 & 0.97 \\
\quad Department of Defense Appropriations Act, 20\dots & 109 & Yea & 0.01 & 0.98 \\
\multicolumn{5}{l}{\emph{Justin Amash (MI)}: 579 held-out votes} \\
\quad Providing for consideration of the Senate am\dots & 115 & Nay & 0.01 & 0.94 \\
\quad Reaffirming the United States' commitment to\dots & 112 & Nay & 0.07 & 0.98 \\
\quad To amend the Ronald Reagan Centennial Commis\dots & 112 & Nay & 0.12 & 0.99 \\
\quad Coastal and Great Lakes Communities Enhancem\dots & 116 & Nay & 0.08 & 0.94 \\
\quad John S. McCain National Defense Authorizatio\dots & 115 & Nay & 0.11 & 0.97 \\
\quad Assisting States' Implementation of Plans of\dots & 115 & Nay & 0.11 & 0.96 \\
\bottomrule
\end{tabular}
}
\end{table}

\begin{figure}[t]
\centering
\includegraphics[width=.68\linewidth]{figures/member_gmp.pdf}
\caption{Member-level fit on identical held-out votes. Each point is
one member (at least 100 held-out votes), plotting held-out geometric
mean probability under this paper's spatial model (vertical axis)
against the frozen-DW-NOMINATE model (horizontal axis); blue denotes
Democrats and red Republicans. Most members lie near the diagonal;
the labeled members, for whom the classical model fits worst, are the
literature's canonical hard cases.}
\label{fig:membergmp}
\end{figure}

\subsection{Forecasting future votes}

Table~\ref{tab:regime} reports the temporal-forecast comparison, in
which every model predicts the final ten percent of each
congress-chamber's rollcalls from the preceding ninety. The ensemble
reaches \champForecastLL{} log loss (GMP \champForecastGMP;
\champForecastAcc\% accuracy, \champContestedAcc\% on contested
votes, APRE \champAPRE) on \forecastTestN{} strictly future votes.
The comparison set is instructive at both ends. A constant-probability
baseline scores \constantForecastLL{}; a member-by-question table of
empirical voting rates scores \memberqForecastLL{}; and the
2011--2016 text-based architectures, reimplemented with modern
optimizers and embeddings where applicable, fail to improve on that
table under proper scoring (Appendix Table~\ref{tab:lit})---the
Gerrish--Blei-style ideal-point-topic model, scored as published,
performs worse than the constant baseline (\gbRawForecastLL{}), and
recovers only to \gbTcalForecastLL{} with temperature calibration.
Those reimplementations do not exhaust what their era's methods might
achieve under indefinite tuning, but the gap between them and the
present pipeline is large. Because the ensemble's TF--IDF component
includes CRS-curated subject metadata whose revision history cannot
be verified (Appendix~\ref{app:audit}), we report throughout both the
ensemble (\champForecastLL) and the fully leakage-clean single
sentence-encoder tower (\incForecastLL); every qualitative conclusion
in this paper holds under the conservative number. Paired
rollcall-cluster bootstraps place the ensemble's margin over its own
strongest component at $[\pairedDeltaOneLo, \pairedDeltaOneHi]$ log
loss, so the gains are not artifacts of vote clustering.

Three stratifications bound what the forecasting result does and does
not mean. First, because bills often generate multiple rollcalls, a
model could in principle recall a bill family's coalition from
training rather than infer it from text. The opposite holds: on the
\freshRC{} held-out rollcalls whose bills were never voted on in the
training window, the text model's advantage over its no-text
counterpart is \emph{larger} (log loss \freshLLchamp{} against
\freshLLnotext) than on the \seenRC{} rollcalls with prior votes on
the same bill (\seenLLchamp{} against \seenLLnotext)---the gains come
from reading new bills, not remembering old ones. Second, the text
gain is concentrated in the House (\houseLLchamp{} against
\houseLLnotext{} without text) and much smaller in the Senate
(\senateLLchamp{} against \senateLLnotext), a chamber difference that
recurs in Section~\ref{sec:protest} and that we report as a boundary
rather than smooth over. Third, the probabilities are usable as
probabilities: expected calibration error is \eceChamp{} for the
ensemble (Brier \brierChamp; reliability curves in
Appendix Figure~\ref{fig:calibration}), against \eceNotext{}
(\brierNotext) for the no-text counterpart.

Because iterative development can overfit even a temporal holdout,
the design includes a test that iteration cannot contaminate: a
predecessor artifact was frozen and hash-pinned on June 12, 2026, and
subsequently scored on the \prospectiveN{} member-vote decisions
Congress cast after its June 9 data snapshot. It recorded
\prospectiveVoneLL{} log loss (\prospectiveVoneAcc\% accuracy) on the
post-snapshot sample and \postFreezeLL{} on the \postFreezeN{}
decisions cast strictly after the freeze date itself, both consistent
with its development-time performance (Table~\ref{tab:prospective}).
The audit is preliminary---\prospectiveRC{} rollcalls is a small
cluster count, and the interval in Table~\ref{tab:prospective} is
correspondingly wide---but it accumulates mechanically as Congress
votes, and the current ensemble is frozen on the same terms.

\begin{table}[t]
\centering
\caption{Temporal-forecast comparison: test log loss and accuracy
under identical evaluation. The right panel reports congress-out
transfer, anticipating Section~\ref{sec:limits}.}
\label{tab:regime}
\begin{tabular}{lcccc}
\toprule
Information used & \multicolumn{2}{c}{Temporal forecast} & \multicolumn{2}{c}{Congress-out} \\
\midrule
Party history only & 0.466 & 76.6 & 0.693 & 57.1 \\
Member history only & 0.459 & 76.0 & 0.693 & 57.1 \\
Rollcall metadata only & 0.511 & 72.3 & 0.717 & 58.6 \\
Text + metadata (no member history) & 0.464 & 77.8 & 0.741 & 64.4 \\
Text + metadata + member history & 0.411 & 80.1 & 0.720 & 59.3 \\
+ modern encoder, deep head, calibration & 0.348 & 85.2 & 0.670 & 61.6 \\
+ ensemble (final) & 0.324 & 86.2 & 0.643 & 62.1 \\
\bottomrule
\end{tabular}

\end{table}

\begin{table}[t]
\centering
\caption{Frozen, hash-pinned artifacts scored on votes cast after
their June 9, 2026 data snapshot (freeze dates June 12 and July 3).
The bootstrap 95\% interval for the v1 log loss is
\prospLLlo--\prospLLhi.}
\label{tab:prospective}
\begin{tabular}{llrrccc}
\toprule
Artifact & Snapshot & Rollcalls & Votes & Log loss & Acc.\ (\%) & AUC \\
\midrule
v1 (emb2-MLP tower) & 2026-06-09 & 43 & 9,451 & 0.331 & 85.7 & 0.924 \\
v2 (three-tower blend) & 2026-06-09 & 43 & 9,451 & 0.332 & 86.8 & 0.921 \\
\bottomrule
\end{tabular}

\end{table}

\section{Text Detects Protest Voting in Advance}
\label{sec:protest}

Protest voting is the scaling literature's newest measurement
problem: majority members voting against bills they arguably prefer
\citep{fowler2026protest}. When such members appear to side with the
opposite pole despite latent preferences near their party's, a
one-dimensional model can misplace them---the ends-against-the-middle
patterns that \citet{duckmayr2023ends} model with a generalized
response function. Existing treatments in both traditions are
retrospective, reinterpreting votes after they are cast. Our
observable is the majority-party \emph{defection}---a broader
category that includes sincere disagreement and constituency pressure
alongside protest---and this section shows that defections are
predictable before the vote, then examines how: which features of the
text carry the signal, and where text alone falls short. The claim
that the predicted defections include the protest episodes of
substantive interest runs through the worked cases.

\subsection{An episode}

The logic has two parts. Member history identifies \emph{who} sits on
a bill's losing flank, but history alone cannot indicate \emph{which
bills} put the flank in revolt, because that depends on what the bill
is---a leadership compromise, a bipartisan spending deal, a
suspension-calendar package. That information is in the text.
Figure~\ref{fig:protest} compares the two models on the American
Relief Act of December 2024 (H.R.~10545), the continuing resolution
that averted a government shutdown and passed 366--34 over a revolt
of the Republican right flank; the vote sits deep in the holdout
window, and neither model observed it. Without text, the architecture
prices the vote as a routine majority-party bill: Democratic support
is uncertain and the revolt is invisible. With the text, the model
recognizes the character of the bill: Democratic probabilities move
to near one, Republican probabilities decline with distance from the
chamber center, and the five members the model considers the
likeliest defectors in the entire chamber---Brecheen, Roy, Massie,
Crane, and Biggs---all voted nay.

\begin{figure}[t]
\centering
\includegraphics[width=.98\linewidth]{figures/protest_detection.pdf}
\caption{One majority-splitting vote (H.R.~10545, December 20, 2024),
predicted by the same architecture without text (left panel) and with
it (right panel). Each point is a member, positioned by first-dimension
ideal point; dots denote members who voted yea (colored by party) and
crosses the 34 members who voted nay. The majority party's modal vote
was yea; ranking all members chamber-wide by predicted probability of
voting nay, the text-reading model's five likeliest---all majority
members, labeled at right---all voted nay.}
\label{fig:protest}
\end{figure}

\subsection{How the model reads the bill}

Where in the text does this information live? We answer with four
diagnostics of increasing directness: what the encoder considers
similar, what a linear probe recovers from the text alone, what
breaks when parts of the text are redacted at scale, and what moves
when specific bills' texts are counterfactually edited.

One clarification first, because it disciplines the language of this
section. The December 2024 continuing resolution appeared the day
before its vote, so no summary existed: the model's entire text input
was the vote question and the title---``On Motion to Suspend the
Rules and Pass. American Relief Act, 2025.'' That is also all a
contemporaneous observer could have read, and it proves sufficient
for Figure~\ref{fig:protest}; but for such late-arriving bills
``reading the bill'' means reading the title and the procedural
posture. Most of the holdout is richer: 75\% of legislation-linked
held-out rollcalls carry a genuine pre-vote summary, and the
diagnostics below use both kinds of case.

The first diagnostic asks what the encoder considers similar.
Table~\ref{tab:neighbors} lists, for three held-out votes of the
118th Congress, the nearest training-window rollcalls in
sentence-embedding space, along with each neighbor's realized
majority-party defection share. The December 2024 continuing
resolution's neighborhood consists of earlier must-pass fiscal
packages---the Tax Relief for American Families and Workers Act, the
bipartisan tax agreement that drew a 22\% majority defection; a
disaster tax-relief bill; and an earlier continuing-appropriations
extension that drew a 42\% defection. The encoder, in other words,
places the bill among precisely the earlier bills that provoked the
same right-flank revolts, and the model's training on those
neighbors' votes is what prices the defection risk. The BIOSECURE
Act, a China-focused measure that passed with near-unanimous majority
support, sits instead among low-defection national-security and
regulatory bills. The second diagnostic removes members from the
picture entirely: a ridge regression of each training rollcall's
majority-defection share on its text embedding alone, evaluated on
the holdout window, predicts realized defection shares at
$r = \revoltProbeR$ across \revoltProbeN{} majority-side passage
votes. Bill text by itself, with no member information, carries a
measurable revolt signal.

The table's middle panel is equally instructive because there text
similarity fails. The Social Security Fairness Act drew a 34\%
majority defection, yet its embedding neighbors are consensual
benefit expansions that drew almost none: nothing in the summary's
language marks the bill as contentious, because the conflict was
procedural---the bill reached the floor by discharge petition rather
than through the majority leadership's agenda. The full model
nevertheless ranks the \ssfaDefectors{} Republican defectors on this
vote better than its no-text twin (within-vote AUC \ssfaAUCtext{}
against \ssfaAUCnotext), drawing on member history and institutional
features rather than content.

The third diagnostic redacts the text at scale.
Table~\ref{tab:ablation} feeds the frozen champion's sentence-encoder
tower systematically corrupted versions of the exact text template it
was trained on, for all \ablRC{} legislation-linked holdout rollcalls,
holding every non-text feature fixed; only the embedding input
changes, so differences isolate the text component. Three results
structure the table. Removing the vote question costs the most
($\ablOrigLL \rightarrow \ablNoQLL$ log loss)---the procedural
posture of a vote is the single most informative phrase the model
reads, consistent with the discharge-petition case above. The placebo
is the sharpest test: replacing a bill's description and summary with
those of a \emph{different} bill in the same policy area and question
bucket costs more ($\ablShuffLL$) than deleting the summary outright
($\ablNoSumLL$)---topic-matched but wrong-bill text actively misleads
the model, so the signal is bill-specific content, not policy-area
membership. And masking proper nouns costs little
($\ablNoPropLL$) while stripping digits costs nothing: the model is
not recognizing names or dates. A policy-area decomposition
(Appendix Table~\ref{tab:policygain}) locates the text gains where
legislative language is most distinctive---government operations,
international affairs, public lands---and near zero for taxation and
commemorative business.

\begin{table}[t]
\centering
\caption{Redaction ablations. Each row re-scores the identical
member-vote cells through the frozen model, changing only the text
supplied to its sentence encoder ($N = \ablRC$ legislation-linked
holdout rollcalls). The defection AUC is the noisier statistic (507
rollcalls) and is shown for completeness; log loss is the scoring
rule.}
\label{tab:ablation}
{\small \begin{tabular}{lccc}
\toprule
Text given to the frozen model & Log loss & $\Delta$ & Defection AUC \\
\midrule
Original text & 0.328 &  & 0.784 \\
Bill summary removed & 0.351 & +0.023 & 0.811 \\
Vote question removed & 0.383 & +0.055 & 0.766 \\
Titles replaced by generic language & 0.340 & +0.012 & 0.768 \\
Proper nouns masked & 0.339 & +0.011 & 0.779 \\
Digits removed & 0.328 & +0.000 & 0.787 \\
Text from another bill, same policy area & 0.392 & +0.064 & 0.777 \\
\bottomrule
\end{tabular}
}
\end{table}

The fourth diagnostic edits individual bills.
Table~\ref{tab:counterfactual} performs counterfactual edits on two
exemplars through the same frozen tower. For the title-only
continuing resolution, replacing the title with generic language cuts
the predicted majority defection from \cfAraOrig\% to \cfAraGeneric\%
---the title's semantics carry real signal---while changing the
question from suspension to ordinary passage \emph{raises} predicted
defection to \cfAraPassage\%, the model having learned that
leadership brings consensus business to the suspension calendar. For
a summary-rich case---the 2006 Appalachian Regional Development
amendments, a regional spending authorization whose 25\% right-flank
revolt the model largely anticipated---substituting a same-area but
different bill's summary halves the predicted revolt
(\cfArdaOrig\% to \cfArdaSwap\%), while deleting interior sentences
that mention spending programs, with the title retained, moves the
prediction almost nowhere (\cfArdaCut\%). Together with the
redaction table, the picture of \emph{how} the model reads is
consistent and honest: it reads the way a hurried professional
does---procedural posture first, then the title and the leading
language of the summary, with the fine print contributing least.

\begin{table}[t]
\centering
\caption{Counterfactual text edits on two exemplar holdout votes,
scored by the frozen model with all non-text features held at their
actual values. ``Top-5 correct'' is the number of the model's five
likeliest majority defectors who in fact voted nay.}
\label{tab:counterfactual}
{\small \begin{tabular}{p{7.6cm}ccc}
\toprule
Text supplied to the frozen model & \shortstack{Predicted\\defection share} & \shortstack{Implied\\defectors} & \shortstack{Top-5 picks\\who defected} \\
\midrule
\multicolumn{4}{l}{\emph{American Relief Act CR, Dec.\ 2024 (pre-vote text is title-only)}} \\
\multicolumn{4}{l}{\quad\emph{realized outcome: 34 of 204 majority members defected (17\%)}} \\
\quad Original (title + suspension question) & 6.5\% & 13 & 5/5 \\
\quad Title replaced by generic language & 2.7\% & 6 & 5/5 \\
\quad Question changed to On Passage & 13.5\% & 28 & 5/5 \\
\quad BIOSECURE Act title substituted & 8.4\% & 17 & 5/5 \\
\multicolumn{4}{l}{\emph{Appalachian Regional Development Act Amendments, 2006 (full pre-vote summary)}} \\
\multicolumn{4}{l}{\quad\emph{realized outcome: 55 of 223 majority members defected (25\%)}} \\
\quad Original (title + summary) & 14.5\% & 32 & 5/5 \\
\quad Spending-program sentences removed & 14.7\% & 33 & 5/5 \\
\quad Generic title, summary kept & 12.6\% & 28 & 4/5 \\
\quad Same-area bill's summary substituted & 7.3\% & 16 & 4/5 \\
\bottomrule
\end{tabular}
}
\end{table}

The mechanism, then, is a division of labor: text locates a bill
among its substantive precedents and sets the coalition's shape, the
procedural posture signals how contested the leadership expects the
vote to be, and member history supplies who sits where---with the
institutional features carrying the cases, like the discharge
petition, in which the path to the floor rather than the language
predicts conflict.

\begin{table}[t]
\centering
\caption{Nearest training-window rollcalls in sentence-embedding
space for three held-out votes of the 118th Congress, with cosine
similarity and each rollcall's realized majority-party defection
share. Neighbors are restricted to votes that precede the exemplar
and lie in the model's training window.}
\label{tab:neighbors}
{\small \begin{tabular}{p{7.8cm}ccc}
\toprule
Nearest earlier bills in the encoder's text space & Congress & \shortstack{Text\\similarity} & \shortstack{Its majority\\defection} \\
\midrule
\multicolumn{4}{l}{\emph{American Relief Act (CR), Dec.\ 2024}: realized majority defection 17\%} \\
\quad Tax Relief for American Families and Workers & 118 & 0.81 & 22\% \\
\quad Federal Disaster Tax Relief Act & 118 & 0.81 & 4\% \\
\quad Further Continuing Appropriations and Other  & 118 & 0.74 & 42\% \\
\multicolumn{4}{l}{\emph{Social Security Fairness Act, Nov.\ 2024}: realized majority defection 34\%} \\
\quad Securing a Strong Retirement Act of 2022 & 117 & 0.69 & 0\% \\
\quad NDO Fairness Act & 118 & 0.68 & 0\% \\
\quad Tax Relief for American Families and Workers & 118 & 0.68 & 22\% \\
\multicolumn{4}{l}{\emph{BIOSECURE Act, Sept.\ 2024}: realized majority defection 1\%} \\
\quad Preventing the Financing of Illegal Syntheti & 118 & 0.70 & 1\% \\
\quad Pro Codes Act & 118 & 0.70 & 26\% \\
\quad Food and Drug Administration Amendments Act  & 110 & 0.69 & 1\% \\
\bottomrule
\end{tabular}
}
\end{table}

\subsection{The result generalizes}

Define the detection task on the \protRC{} held-out rollcalls in
which a member's party majority voted yea and at least three majority
members defected (\protN{} defections in total): rank that party's
members by predicted defection probability, compute the AUC within
each rollcall, and average with equal weight across rollcalls. The
text-reading ensemble attains \protAUCtext{}, against \protAUCnotext{}
for its no-text counterpart---text resolves roughly a fifth of the
remaining ranking errors. The result is robust in the directions a
skeptic would probe. The fully leakage-clean single tower attains
\protAUCclean, so the gain does not depend on the CRS-curated
metadata discussed in Appendix~\ref{app:audit}. Dropping the
three-defection filter and averaging over all \protRCall{}
positive-defection rollcalls gives \protAUCtextAll{} against
\protAUCnotextAll. Predicted defection \emph{shares} track realized
shares across rollcalls (predicted--realized correlation
\defShareRChamp{} against \defShareRNotext{} without text), though
the highest-decile predictions overshoot---the model flags the right
votes but overstates how large the largest revolts will be
(Appendix Figure~\ref{fig:calibration}b). The one boundary is the
chamber split: the text gain lives in the House
(\protAUCtextHouse{} against \protAUCnotextHouse) and is absent in
the Senate (\protAUCtextSenate{} against \protAUCnotextSenate{} over
\protRCsenate{} rollcalls), where the negotiated floor agenda gives
bill text little marginal purchase---a finding we report as a limit
of the instrument, not a footnote.

Two connections tie the result to the measurement literature, one of
them chastened by its own check. The canonical hard cases of
Figure~\ref{fig:membergmp} are also members whose defections the
text model prices well---Paul and Amash sit at the
\paulDefGainPct th and \amashDefGainPct th percentiles of the
member-level text gain---but a member-level
overlap test shows the two improvements are essentially uncorrelated
across the full membership (Spearman $\rho = \overlapRho$ over
\overlapN{} members): the scaling gains and the defection gains are
related at the level of the phenomenon, not member by member, and we
claim no more. And the ensemble's largest remaining
residuals---defections
inconsistent with even the full behavioral record---are released as a
ranked ledger, a sampling frame enriched for pressure and defection
episodes, for studying the votes that even text cannot explain.

\section{Text Predicts Bill Cutpoints}
\label{sec:cutpoints}

The bill-location literature recovers proposal and status-quo
positions from realized votes and cosponsorships
\citep{peress2013estimating, richman2011parties}---necessarily after
the fact. If bill content organizes coalitions, however, a bill's
text should locate it before anyone votes. For every rollcall we take
the realized cutpoint and direction from the spatial fits, hold out
the final twenty percent of each congress-chamber's rollcalls, and
predict both quantities from strictly pre-vote features.

Figure~\ref{fig:cutpred} and Table~\ref{tab:cutpred} report the
result. On the \cutTestN{} identified held-out rollcalls, text plus
metadata predicts cutpoint locations with a mean absolute error of
\cutComboMAE{} member-standard-deviations and a predicted-realized
correlation of $r = \cutComboR$, against a no-feature baseline of
\cutConstMAE{} (and $r = 0$); it classifies coalition direction
correctly for \dirCombo\% of rollcalls, against a
majority-direction baseline of \dirConst\%. The headline feature set
here contains only the sentence-embedding representation of pre-vote
text plus institutional metadata; the CRS-revisable subject fields
discussed in Appendix~\ref{app:audit} enter only the separately
reported TF--IDF row. Panel~(a) shows the predictions
member-standard-deviation by member-standard-deviation: mass
concentrates along the diagonal, with the labeled landmark bills
indicating the scale's interpretation---bills realized to cut deep in
the majority's territory are predicted there.

\begin{figure}[t]
\centering
\includegraphics[width=.98\linewidth]{figures/cutpoint_prediction.pdf}
\caption{Predicting held-out rollcalls' cutpoints from pre-vote
features. Panel (a): predicted (horizontal) against realized
(vertical) cutpoints under the text-plus-metadata model, in
member-standard-deviation units, with selected passage votes labeled;
the gray line is the diagonal. Panel (b): mean absolute error of the
cutpoint prediction by feature set; the dotted line is the no-feature
baseline. Panel (c): share of rollcalls whose coalition direction is
classified correctly; the dotted line is the majority-direction
baseline.}
\label{fig:cutpred}
\end{figure}

\begin{table}[t]
\centering
\caption{Predicting held-out rollcalls' cutpoints and directions from
pre-vote features. The bottom panel decomposes the text into its
components.}
\label{tab:cutpred}
\begin{tabular}{lccc}
\toprule
Features & Cutpoint MAE & Cutpoint $r$ & Direction acc.\ (\%) \\
\midrule
Constant & 0.856 & 0.00 & 55.8 \\
Metadata only & 0.657 & 0.48 & 70.8 \\
Text: TF--IDF & 0.742 & 0.43 & 61.3 \\
Text: embeddings & 0.735 & 0.42 & 59.5 \\
Text + metadata & 0.637 & 0.55 & 74.5 \\
\quad question only & 0.763 & 0.41 & 58.2 \\
\quad description only & 0.813 & 0.25 & 57.2 \\
\quad bill summary only & 0.765 & 0.38 & 59.7 \\
\bottomrule
\end{tabular}

\end{table}

The decomposition in panels (b) and (c) assigns each input its role.
Institutional metadata---question type, the bill sponsor's party,
and, on amendment votes, the amender's party, with their
interactions---is a strong baseline on its own: it locates cutpoints
with $\mathrm{MAE} = \cutMetaMAE$ and supplies a default coalition
side correct \dirMeta\% of the time, which is what one would expect
in a chamber whose agenda is set by the majority party. Text improves
every metric beyond that baseline, with the largest gains in
direction (\dirCombo\%) and location correlation ($\cutMetaR
\rightarrow \cutComboR$). The interpretation is the same division of
labor as in Section~\ref{sec:protest}: the acting legislator's party
indicates which side a bill is probably on, and the content indicates
when that default is wrong, and by how much. Among the text
components (Table~\ref{tab:cutpred}, bottom panel), the bill summary
carries most of the directional information and the vote question
most of the locational information.

Figure~\ref{fig:dirterms} examines the directional signal at the term
level, using an interpretable sparse model fit to final-passage
summaries only. The language of program-building---\emph{establish},
\emph{grants}, \emph{assistance}, \emph{workers},
\emph{emergency}---is associated with liberal-yea coalitions, and the
language of restraint and revision---\emph{rule}, \emph{regulation},
\emph{permit}, \emph{budget}, \emph{injury}---with conservative-yea
ones. These are predictive associations on a leader-curated agenda,
not causal claims about language; but they indicate that the
embedding models' directional signal rests on substantively sensible
regularities rather than incidental correlates.

Because the prediction target is itself an estimate, two
stratifications gauge how much of the residual error belongs to the
target rather than the prediction
(Appendix Table~\ref{tab:cutstrata}). Realized cutpoints are pinned
most precisely where rollcall discrimination is strong; on the
high-discrimination half of the holdout, the mean absolute error
falls to \cutHiDiscMAE{} member-standard-deviations and the
predicted--realized correlation rises to $r = \cutHiDiscR$, so a
substantial share of the headline error is noise in the target. By
vote type, resolutions and procedural votes are located most
accurately and cloture least; final-passage votes sit at
$r = \cutPassageR$. Two caveats bound the
exercise. The prediction target is the cutpoint of our own spatial
fit---a realized-vote quantity with estimation error of its own---and
we have not yet benchmarked the predictions against the
proposal/status-quo estimates of \citet{peress2013estimating} or
\citet{richman2011parties} on an overlapping sample; that comparison,
which requires reimplementing those estimators, is the natural next
test. An interactive forecaster operationalizes the section's result:
a user pastes a hypothetical bill's text, selects a sponsoring party
and chamber, and receives a predicted cutpoint and a member-by-member
vote forecast, with a sponsor--coalition mismatch flag for inputs off
the curated agenda (Appendix~\ref{app:dashboard}).

\begin{figure}[t]
\centering
\includegraphics[width=.62\linewidth]{figures/direction_terms.pdf}
\caption{Bill-summary terms most associated with coalition direction,
from a sparse linear model on final-passage votes. Positive
coefficients (red) indicate language associated with conservative-yea
coalitions; negative coefficients (blue), liberal-yea coalitions.}
\label{fig:dirterms}
\end{figure}

\section{Two Boundaries}
\label{sec:limits}

Users of these methods should know where they stop, and both
boundaries are measured rather than conjectured. The first concerns
\emph{majority transitions}. At every House majority flip in the
window---and at none of the placebo transitions---the
procedural-vote component of member history stops transferring: the
model's procedural log loss averages \flipProcLL{} across flips
against \placeboProcLL{} across placebos, while final-passage
transfer is flip-invariant (\flipPassLL{} against \placeboPassLL;
Appendix Figure~\ref{fig:transitions}). Much of what any vote-based
measure captures on the procedural agenda is institutional role
rather than stable preference, and scores estimated in one majority
configuration should be carried across configurations only with care
\citep{bateman2017house}. The second concerns \emph{amendments}.
Supplying amendments' purpose text makes prediction worse, not
better, in pre-registered, placebo-controlled experiments
(Appendix~\ref{app:ledger}); what moves amendment prediction is actor
identity---the amender's party takes direction accuracy from
approximately a coin flip to 77\%. Short purpose text adds nothing
once bill identity and procedural context are modeled, a measured
limit on what legislative language reveals, at least in the
representations tested here.

\section{Conclusion}

Reading the bills makes Congress more measurable in three specific
ways: a fit standard that improves on the classical measure exactly
where it struggles; advance detection of the protest voting that
scaling models could previously only reinterpret, with
embedding-space evidence of how the text carries the signal; and
prospective bill locations that the agenda literature could
previously estimate only in retrospect. The benchmark, the frozen
artifacts whose forecasting claims the future continues to audit, the
interactive forecaster, and a complete ledger of what failed along
the way are released together.

\bibliographystyle{apsr}
\bibliography{references}

\appendix

\section{Additional Tables}
\label{app:tables}

\begin{table}[h]
\centering
\caption{Completion regime (held-out vote completion; same-rollcall
votes observable): the psychometric ladder.}
\label{tab:regimeA}
\begin{tabular}{lcccc}
\toprule
Model & Log loss & Acc.\ (\%) & Contested acc.\ (\%) & APRE \\
\midrule
Constant rate & 0.655 & 63.7 & 50.9 & -0.32 \\
Member rate & 0.629 & 64.9 & 55.2 & -0.28 \\
Party $\times$ rollcall majority & 0.206 & 91.7 & 92.8 & 0.70 \\
Ideal points (1D) & 0.193 & 91.6 & 91.2 & 0.69 \\
DW-NOMINATE logit (frozen scores) & 0.164 & 93.3 & 93.9 & 0.76 \\
Ideal points (2D) & 0.161 & 93.3 & 93.8 & 0.76 \\
Ideal points (8D) & 0.153 & 93.8 & 94.3 & 0.77 \\
\bottomrule
\end{tabular}

\end{table}

\begin{table}[h]
\centering
\caption{The 2011--2016 text-based architectures, reimplemented and
evaluated under identical conditions.}
\label{tab:lit}
\begin{tabular}{lccccc}
\toprule
Model & Fcst.\ LL & Fcst.\ AUC & Fcst.\ cont.\ acc.\ (\%) & Rand.\ LL & Rand.\ AUC \\
\midrule
Constant rate & 0.617 & 0.500 & 52.7 & 0.659 & 0.500 \\
Party $\times$ question rate & 0.466 & 0.805 & 69.8 & 0.504 & 0.797 \\
Member $\times$ question rate & 0.459 & 0.812 & 69.2 & 0.494 & 0.807 \\
NOMINATE $\times$ context logit & 0.532 & 0.740 & 59.6 & 0.587 & 0.713 \\
GB-style spatial (TF-IDF) & 0.803 & 0.689 & 57.6 & 1.054 & 0.687 \\
GB-style spatial (TF-IDF, calibrated) & 0.566 & 0.689 & 57.1 & 0.685 & 0.687 \\
GB-style spatial (mod.\ emb., calibrated) & 0.556 & 0.699 & 55.8 & 0.662 & 0.670 \\
Kraft-style bilinear & 0.497 & 0.790 & 67.2 & 0.522 & 0.794 \\
TF-IDF tower + member offset & 0.411 & 0.860 & 75.0 & 0.448 & 0.854 \\
MiniLM-MLP tower (prior champion) & 0.348 & 0.906 & 84.0 & 0.417 & 0.880 \\
Three-tower blend (champion) & 0.324 & 0.919 & 85.1 & 0.390 & 0.896 \\
\bottomrule
\end{tabular}

\end{table}

\begin{figure}[h]
\centering
\includegraphics[width=.95\linewidth]{figures/transitions.pdf}
\caption{Majority-transition study: cross-congress transfer at every
House transition (red: majority flips), overall and by vote type.}
\label{fig:transitions}
\end{figure}

\begin{table}[h]
\centering
\caption{The text gain (no-text minus model log loss, positive =
text helps) by CRS policy area, passage-family holdout votes; areas
with at least 25 held-out rollcalls. Policy area is an evaluation
stratifier only, never a model input.}
\label{tab:policygain}
{\small \begin{tabular}{lccc}
\toprule
Policy area & Held-out rollcalls & \multicolumn{2}{c}{Text gain in log loss} \\
 & & Ensemble & Clean tower \\
\midrule
Government Operations and Politics & 64 & +0.103 & +0.094 \\
International Affairs & 40 & +0.090 & +0.075 \\
Public Lands and Natural Resources & 46 & +0.075 & +0.062 \\
Armed Forces and National Security & 90 & +0.051 & +0.059 \\
Transportation and Public Works & 26 & +0.044 & +0.035 \\
Economics and Public Finance & 51 & +0.012 & +0.017 \\
Crime and Law Enforcement & 58 & +0.058 & +0.011 \\
Energy & 28 & +0.044 & +0.005 \\
Health & 50 & +0.052 & -0.006 \\
Taxation & 43 & +0.042 & -0.021 \\
Congress & 165 & -0.006 & -0.028 \\
Finance and Financial Sector & 43 & +0.022 & -0.070 \\
\bottomrule
\end{tabular}
}
\end{table}

\begin{table}[h]
\centering
\caption{Cutpoint prediction (text + metadata) by vote type and by
target precision. Realized cutpoints are estimated quantities;
strongly discriminating rollcalls pin them most precisely.}
\label{tab:cutstrata}
{\small \begin{tabular}{lccc}
\toprule
Subset & Rollcalls & MAE & $r$ \\
\midrule
All rollcalls & 3,527 & 0.619 & 0.55 \\
Passage & 666 & 0.821 & 0.61 \\
Amendment & 929 & 0.516 & 0.49 \\
Procedural & 670 & 0.410 & 0.49 \\
Cloture & 381 & 0.744 & 0.43 \\
Resolution & 319 & 0.299 & 0.70 \\
High discrimination ($|a_v|$ above median) & 1,763 & 0.352 & 0.70 \\
Low discrimination ($|a_v|$ below median) & 1,764 & 0.884 & 0.51 \\
\bottomrule
\end{tabular}
}
\end{table}

\begin{figure}[h]
\centering
\includegraphics[width=.95\linewidth]{figures/calibration.pdf}
\caption{Calibration on the temporal-forecast holdout. Panel (a):
member-vote reliability curves (16 equal-count bins). Panel (b):
predicted against realized majority-defection shares by decile of the
prediction, majority-yea passage-family rollcalls with at least 50
majority members.}
\label{fig:calibration}
\end{figure}

\section{Measurement Provenance}
\label{app:provenance}

Which estimator produces each quantity:
the ensemble produces vote forecasts (pre-vote), the protest-detection
scores (pre-vote), and the out-of-character residual ledger
(in-sample, labeled as such); per chamber-congress spatial fits
produce positions, cutpoints, and directions (from realized votes);
ridge and logistic models on pre-vote features produce the cutpoint
\emph{predictions} of Section~\ref{sec:cutpoints}. An identifiability
remark bounds the enterprise: if two member-rollcall states are
observationally equivalent given the full available information
set---inducing the same joint distribution of votes and observed
covariates---no measure computed from that information can
distinguish them. The
repository's synthetic suite demonstrates the constructive direction:
planted structures are recovered exactly by the models whose held-out
scores improve.

\section{Feature Timing and Text Coverage}
\label{app:audit}

Every input and its timestamp rule: vote question and description
(fixed at the vote); bill title and sponsor (introduction); bill
summaries (only versions dated strictly before the vote); amendment
purpose and amender party (filed at submission; joined through
recorded-vote references with session-offset correction);
member-by-question rates (training window only); majority status
(member composition); early stopping and calibration (final 5\% of
training rollcalls, inside train). CRS-revisable policy metadata is
excluded from the leakage-clean text; the ensemble's TF--IDF
component includes it, so the fully conservative headline is the
single MiniLM tower's \incForecastLL{} rather than the ensemble's
\champForecastLL. Text coverage: the vote question is universal; 95\%
of legislation-linked rollcalls have a pre-vote summary or title;
\amdtFilledNoDesc{} of \amdtNoDesc{} text-starved amendment rollcalls
gain amendment-side information via the records join.

\section{Model Specification}
\label{app:spec}

The champion is a logistic blend (weights 0.57/0.12/0.42, fit on the
temporal development slice) of three two-tower models: member
embedding $x_i \in \mathbb{R}^{16}$ and intercept; rollcall
parameters amortized from features by a two-layer MLP (width 128);
sponsor-alignment and member-by-question logit offsets; temperature
calibration on the development slice. Towers differ in text features:
MiniLM (all-MiniLM-L6-v2, frozen, 384d, 1,500-character template);
plus TF--IDF/SVD of title-subject metadata; and Qwen3-Embedding-0.6B
(1,024d, 6,000 characters). AdamW (lr .02, wd $10^{-4}$), batch
$2^{17}$, early stopping on the temporal slice. The blend weights are
unconstrained logistic-stacking coefficients (they need not sum to
one). The no-text counterpart deletes all text blocks and is
independently trained, early-stopped, and temperature-calibrated
under the identical protocol; nothing is inherited from the text
models. One
pipeline command regenerates every exhibit from raw Voteview and
BILLSTATUS files; the headline table costs under two GPU-hours on a
laptop-class accelerator.

\section{The Experiment Ledger}
\label{app:ledger}

Every development experiment was pre-registered with hypothesis,
knowability argument, and falsification test before its first run;
selection used validation only; the final three-model comparison
observed test sets once. Recorded negatives, with mechanisms: encoder
scale (a $16\times$ text budget does not help and worsens long
bills); amendment purpose text (three variants and a placebo control;
Section~\ref{sec:limits}); recency-weighted member rates; within-bill
context features; per-bucket calibration; and, recorded within the
E4b entry, regime inheritance (parameters fit on the temporal
development slice transfer to temporal evaluation and degrade on
interleaved evaluation).

\begin{table}[h]
\centering
\caption{Ledger audit trail (full record with timestamps in the
repository).}
\label{tab:ledger}
{\small
\begin{tabular}{llll}
\toprule
ID & Hypothesis & Falsification test & Outcome \\
\midrule
M-A & bigger encoder helps long bills & long-bill val loss falls & negative \\
E1 & within-bill context helps & beats shuffled-bill placebo & negative \\
E1b & stacked context corrector & permuted-bill weights $\to 0$ & null \\
E2 & recency-weighted rates help & $\tau\to\infty$ recovers flat & negative \\
E3 & per-bucket calibration helps & val loss below global cal. & negative \\
E4/E4b & tower blends help & dev-weight collapse if redundant & positive \\
E5a--c & amendment text helps & beats placebo file rebuild & negative \\
E4b(r) & dev-slice params inherit regime & interleaved val degrades & confirmed \\
S5 & blend beats champion & one-shot test LL below champion & positive \\
\bottomrule
\end{tabular}}
\end{table}

\section{The Interactive Forecaster}
\label{app:dashboard}

A single-tower approximation of the frozen champion (MiniLM component
only; near-threshold forecasts should be read accordingly). Pasted
text enters the training template and encoder; a synthetic bill
record carries the user's sponsoring party; the cutpoint comes from a
fractional-response logistic fit of predicted probabilities on
current-member positions, with fitted mean $q_i =
\sigma(\alpha + \beta x_i)$; the reported cutpoint $x^* =
-\alpha/\beta$ is the 0.5 crossing of the fitted curve, displayed
only when the slope is meaningful and the crossing lies within the
chamber. Forecasts whose predicted coalition conflicts with the
bill's own sponsorship carry a sponsor--coalition mismatch flag: such
inputs lie off the curated floor agenda on which the model was
trained, and its extrapolations there are unvalidated.

\end{document}
